\documentclass[11pt,a4paper]{article}
\usepackage[margin=1in]{geometry}
\usepackage{amsmath,amssymb}
\usepackage{booktabs}
\usepackage{hyperref}
\usepackage{xcolor}
\usepackage{float}

\definecolor{good}{rgb}{0.0,0.5,0.0}
\definecolor{bad}{rgb}{0.7,0.0,0.0}
\definecolor{improved}{rgb}{0.0,0.3,0.7}

\title{EEG Scalp-to-In-Ear Signal Prediction:\\
  Iteration 003 --- Scaling Data and Training}
\author{Automated Research Loop}
\date{April 7, 2026}

\begin{document}
\maketitle

\begin{abstract}
We scaled the synthetic training data from 5 to 20 subjects (4$\times$ more data),
reduced the FIR filter length from 65 to 17 taps, and increased Conv Encoder
training from 300 to 500 epochs. The Conv Encoder achieves $r=0.876$
(SNR$=6.33$\,dB), now meeting convergence criteria ($r \geq 0.85$). The FIR
filter remains stalled at $r=0.695$, suggesting a fundamental architectural
mismatch with this purely spatial problem. Three of four models now converge;
Iteration~004 will address the FIR filter with closed-form initialization of
its center tap.
\end{abstract}

\section{Changes from Iteration 002}

\begin{enumerate}
  \item \textbf{4$\times$ more training data}: \texttt{n\_subjects}: $5 \to 20$,
    \texttt{n\_samples}: $76{,}800 \to 153{,}600$, yielding $\sim$8,360 training
    windows (vs.\ 2,090 previously).
  \item \textbf{Shorter FIR filter}: $L = 65 \to 17$ taps ($\sim$66\,ms context),
    reducing parameters from 5,460 to 1,428. Hypothesis: the true forward model
    is purely spatial (zero-lag), so fewer temporal parameters should reduce
    overfitting.
  \item \textbf{More Conv Encoder training}: $300 \to 500$ epochs with cosine
    annealing $T_{\max} = 500$.
\end{enumerate}

\section{Results}

\begin{table}[H]
\centering
\caption{Iteration 003 vs 002 comparison.}
\begin{tabular}{lrrrrrr}
\toprule
& \multicolumn{2}{c}{Pearson $r$} & \multicolumn{2}{c}{SNR (dB)} & \multicolumn{2}{c}{Coherence} \\
\cmidrule(lr){2-3} \cmidrule(lr){4-5} \cmidrule(lr){6-7}
Model & It.~2 & It.~3 & It.~2 & It.~3 & It.~2 & It.~3 \\
\midrule
Closed-form & 0.887 & 0.887 & 6.51 & 6.51 & 0.715 & 0.715 \\
Linear SGD & 0.887 & 0.887 & 6.51 & 6.51 & 0.715 & 0.715 \\
FIR Filter & 0.695 & 0.695 & 2.86 & 2.87 & 0.333 & 0.334 \\
Conv Encoder & 0.815 & \textcolor{good}{0.876} & 4.75 & \textcolor{good}{6.33} & 0.442 & \textcolor{good}{0.552} \\
\bottomrule
\end{tabular}
\end{table}

\subsection{Analysis}

\paragraph{Conv Encoder breakthrough.}
The combination of 4$\times$ more data and extended training (500 epochs) boosted
the Conv Encoder from $r=0.815$ to $r=0.876$, a $\Delta r = +0.061$ improvement.
SNR jumped from 4.75 to 6.33\,dB, approaching the closed-form's 6.51\,dB.
The remaining gap ($r^2 = 0.767$ vs $0.787$) is only 2\% of variance, likely
attributable to the Conv Encoder's nonlinear capacity being slightly harmful
for this linear problem.

\paragraph{FIR filter plateau.}
Reducing filter length from 65 to 17 taps had negligible effect ($\Delta r < 0.001$).
The FIR architecture applies a temporal convolution \emph{per input channel},
then sums across channels. Even with $L=17$, it has 1,428 parameters learning
temporal filters where the true signal is purely spatial. The temporal freedom
allows it to fit noise correlations rather than the spatial mixing.

\paragraph{Diagnosis: why FIR lags.}
The FIR filter's 1D convolution learns $\mathbf{W}[k] \in \mathbb{R}^{4 \times 21}$
for each lag $k \in \{0, \ldots, L-1\}$. For the true model $\mathbf{Y}[t] =
\mathbf{M}\mathbf{X}[t]$, only $\mathbf{W}[0] = \mathbf{M}$ matters; all other
lags should be zero. But SGD distributes energy across taps, fitting temporal
noise autocorrelations. The fix: initialize $\mathbf{W}[\lfloor L/2 \rfloor]$
from $\mathbf{W}^*$ and all other taps to zero.

\section{Convergence Status}

\begin{table}[H]
\centering
\begin{tabular}{lcccc}
\toprule
Criterion & Closed-form & Linear SGD & FIR & Conv \\
\midrule
$r \geq 0.85$ & \textcolor{good}{\checkmark} & \textcolor{good}{\checkmark} & \textcolor{bad}{$\times$} (0.695) & \textcolor{good}{\checkmark} \\
SNR $\geq 5.0$\,dB & \textcolor{good}{\checkmark} & \textcolor{good}{\checkmark} & \textcolor{bad}{$\times$} (2.87) & \textcolor{good}{\checkmark} \\
No negative $r$ & \textcolor{good}{\checkmark} & \textcolor{good}{\checkmark} & \textcolor{good}{\checkmark} & \textcolor{good}{\checkmark} \\
\bottomrule
\end{tabular}
\end{table}

\textbf{3/4 models converged.} FIR filter is the sole remaining bottleneck.

\section{Plan for Iteration 004}

\begin{enumerate}
  \item \textbf{Closed-form FIR initialization}: Set the center tap
    $\mathbf{W}[\lfloor L/2 \rfloor]$ to $\mathbf{W}^*$ and all other taps to
    zero. This gives the FIR a strong starting point.
  \item \textbf{Regularize non-center taps}: Add L1 penalty on non-center taps
    to encourage the model to keep energy concentrated at zero lag.
  \item If FIR still underperforms, accept that the architecture is inappropriate
    for purely spatial problems and declare convergence for the other three models.
\end{enumerate}

\end{document}
